% Thorne & Campolattaro 1967, ApJ 149, 591
% Page 9 (p. 599) transcription

solution for which the Lagrangian change in the pressure does not vanish. For the
example of a polytropic pressure-density relation (eq.~[17a]) this unacceptable solution has
fluid displacements which diverge at the star's surface ($W \sim V \sim [R - r]^{-8}$); see
Appendix~D for details.

For any arbitrary choice of the complex pulsation frequency, $\omega$, there will be just one
solution to the eigenequations~(14) which is well-behaved both at the center of the star
and at the surface. That solution is the unique linear mixture of the two solutions
well behaved at the center, which contains none of the solution unacceptable at the
surface. Henceforth we confine our attention to that unique solution, for any given
choice of $\omega$, which is acceptable both at the center and surface.

\subsection*{c) Behavior of the Physically Acceptable Eigensolution at $r = \infty$}

For any given choice of $\omega$ the physically acceptable eigensolution will be
characterized by the behavior it takes on outside the star by some particular mixture of
ingoing and outgoing gravitational radiation. As is shown in Appendix~E, the gravitational
waves have the asymptotic form (solution to eigenqs.~[14a] for large $r$):
\begin{align}
K &= C^{(\mathrm{I})} g_l^{(+)}(r)\, e^{-i(\omega t + \text{SM}\ln r)} + C^{(\mathrm{O})} g_l^{(-)}(r)\, e^{-i(\omega t - \text{SM}\ln r)} \notag \\
H &= i a C^{(\mathrm{I})} g_l^{(+)}(r)\, e^{-i(\omega t + \text{SM}\ln r)} - i a C^{(\mathrm{O})} g_l^{(-)}(r)\, e^{-i(\omega t - \text{SM}\ln r)} \notag \\
&\quad \text{for } r \gg M \text{ and } \{R \to r\}^{-1} \notag
\end{align}
\begin{equation}
\label{eq:18}
\tag{18}
\end{equation}
Here $C^{(\mathrm{I})}$ and $C^{(\mathrm{O})}$ are complex constants which are fixed by the demand that the
external perturbation in the geometry of spacetime join smoothly at the star's surface
to the internal perturbation

\begin{equation}
K,\; K',\; H \text{ and } H \text{ continuous at } r = R \;.
\label{eq:19}
\tag{19}
\end{equation}

Since the time dependence of the normal modes is $e^{i\omega t}$ (cf.\ eq.~[12]), \textit{the constant $C^{(\mathrm{I})}$
is the amplitude for incoming gravitational waves, while $C^{(\mathrm{O})}$ is the amplitude for outgoing
gravitational waves}. For any given choice of the complex eigenfrequency, $\omega$, and of the
spherical harmonic index, $l$, one has a uniquely determined ratio of outgoing amplitude
to incoming amplitude:

\begin{equation}
C^{(\mathrm{O})}/C^{(\mathrm{I})} \text{ is a unique function of } \omega \text{ and } l \;.
\label{eq:20}
\tag{20}
\end{equation}

From previous experience with one-dimensional analyses one expects this---that, for
any given choice of $l$, there should be a discrete spectrum of complex eigenfrequencies
for which $C^{(\mathrm{O})}/C^{(\mathrm{I})}$ is infinite. These eigenfrequencies and the corresponding eigenfunctions make up normal modes with purely outgoing radiation. Similarly, those discrete
eigenfrequencies and functions for which $C^{(\mathrm{O})}/C^{(\mathrm{I})}$ is zero are normal modes with
purely incoming radiation. Finally, those normal modes with real eigenfrequencies
($r = \omega$ in eq.~[11]) have real eigenfunctions (cf.\ eqs.~[14]) and represent physically
acceptable, standing waves with equal inward and outward fluxes.

Most of the eigenfunctions~(18) far from the star might appear at first sight to be
physically unacceptable. In general, they contain a divergence in the metric perturbation
tensor (eqs.~[5], [7b], and [12]) which diverges at large $r$ rather than having the $1/r$
behavior that one expects of a radiation field. The divergence of the metric perturbation
at large $r$ arises from two sources:

1. There is an exponential divergence associated with the imaginary part, $i/r$, of $\omega$.
This exponential divergence is physically acceptable. For $r > 0$ the divergence occurs in the outgoing-wave term and is a result of the exponential damping of stable pulsations of the star.
For $r < 0$ the divergence occurs in the incoming-wave term and is associated with an
input of energy from infinity which increases exponentially with time.

2. There is also a divergence of order $r$ in the amplitude of the metric perturbation
